\documentclass[10pt,a4paper]{amsart}
\usepackage[margin=19mm]{geometry}
\usepackage{amsmath,amssymb,mathtools}
\usepackage[scr=rsfs,cal=euler]{mathalfa}
\usepackage{xfrac}
\usepackage[unicode,hidelinks,bookmarks=false]{hyperref}
\newcommand{\ie}{i.e.\ }
\newcommand{\cf}{cf.\ }
\newcommand{\resp}{respectively\ }
\newcommand{\Subsection}{\subsection}
\providecommand{\nobreakdash}{\nobreak\hbox{-}}
\setlength{\emergencystretch}{2em}
\multlinegap0pt
\usepackage{mathrsfs}
\usepackage{booktabs}
\usepackage{textcomp}
\usepackage{enumitem}
\usepackage{appendix}
\usepackage{url}
\newcommand{\R}{\mathbb{R}}
\newcommand{\CC}{\mathbb{C}}
\newcommand{\eps}{\varepsilon}
\newcommand{\dd}{\,d}
\DeclareMathOperator{\dist}{dist}
\newcommand{\backmatter}{}
\newcommand{\bmhead}[1]{\section*{#1}}
\newtheorem{theorem}{Theorem}
\newtheorem{proposition}[theorem]{Proposition}
\newtheorem{lemma}[theorem]{Lemma}
\newtheorem{corollary}[theorem]{Corollary}
\newtheorem{Assumption}{Assumption}
\newtheorem{assumption}[theorem]{Assumption}
\newtheorem{Lemma}{Lemma}
\newtheorem{Corollary}{Corollary}
\newtheorem{example}{Example}
\newtheorem{remark}{Remark}
\newtheorem{definition}{Definition}
\begin{document}
\title[A Priori Integral Persistent Excitation in Conservative Poly]{A Priori Integral Persistent Excitation in Conservative Polynomial ODEs with Higher-Order Interactions}
\author{Aleksandr Semenov; Alexander Fradkov}
\date{}
\maketitle
\noindent\emph{Source and licence.} A Priori Integral Persistent Excitation in Conservative Polynomial ODEs with Higher-Order Interactions. Version arXiv:2607.00091v1 (CC BY 4.0). This material is distributed under the Creative Commons Attribution 4.0 International licence, http://creativecommons.org/licenses/by/4.0/, as recorded in the arXiv OAI licence metadata for this exact version and shown on the abstract page https://arxiv.org/abs/2607.00091v1. Full rights notice: licenses/arxiv-2607.00091-CC-BY-4.0.txt.\par\smallskip
\noindent\textit{Editorial scope.} Counted unit: the original \texttt{theorem} environment \texttt{thm:main\_nonsticking\_2} at source lines 211--219 together with its complete proof at lines 220--222; the delivered document keeps source lines 137--231 so that every referenced label and every citation resolves inside it. Native editable source is the paper's own concatenated TeX; the AMS wrapper is editorial formatting only. No publisher class, upstream script or unrelated artwork is redistributed.
\begin{equation}\label{eq:polynomial_system_2}
    \dot x=f_\xi(x),\qquad x\in\R^n,
    \quad \xi\in\Xi\subset\R^m,
\end{equation}
where the components of $f_\xi$ are polynomials in $x$ and the parameters enter the right‑hand side linearly:
\begin{equation}\label{eq:linear_parameterization_2}
    \dot{x}_i=\sum_{j}\xi_{ij}p_{ij}(x).
\end{equation}
Here $p_{ij}(x)$ are monomials. 

This formulation covers many models of interest, including generalized Lotka--Volterra system, oscillator networks, resonant cascades, and other complex systems.


In the sequel, we will drop the subscript $\xi$ wherever this causes no confusion and it is clear from the context which parameter value is meant; in particular, we will write $f$, $K$, $\mu$, $\Phi^t$, $M$, and $Q$ instead of $f_\xi$, $K_\xi$, $\mu_\xi$, $\Phi_\xi^t$, $M_\xi$, and $Q_\xi$.

\begin{assumption}[Conditions]
\label{ass:main_2}
Let $K\subset\R^n$ be a compact invariant set of system \eqref{eq:polynomial_system_2}. Assume that the following conditions hold.
\begin{enumerate}[label=(A\arabic*)]
    \item $\Phi^t(K) \subset K$ for all $t\in\R$.
    \item On $K$ there exists an invariant probability measure $\mu$ that is absolutely continuous with respect to the Lebesgue measure, and
    \[
        d\mu=\rho(x)\dd x,
        \qquad
        0<c\le \rho(x)\le C<\infty
        \quad \text{on }K.
    \]
    %\item There are no stationary points on $K$: $f(x)\ne0$ for all $x\in K$.
    \item For the matrix
    \[
        M(x)=\bigl[f(x)\mid L_f f(x)\mid\dots\mid L_f^{n-1}f(x)\bigr]
    \]
    the polynomial
    \[
        Q(x)=\det M(x) \not\equiv 0
    \]
    is nondegenerate: there exists a point $x_*\in K$ such that $Q(x_*)\ne0$.
\end{enumerate}
\end{assumption}

Notice that this list of conditions is not overly restrictive. In particular, any Hamiltonian system satisfies the first two conditions. Indeed, let $W(x)$ be the Hamiltonian of the system; then the set $K = \{x:c_1\le W(x) \le c_2\}$, where $c_1,c_2 > 0$, is compact and invariant. By Liouville's theorem, a Hamiltonian system preserves the Lebesgue measure. The third condition is technical and can usually be guaranteed for almost all parameter values. For an easy check, the following statement can be applied.

\begin{proposition}[Checking nondegeneracy for almost all parameters]
\label{prop:param_nondegeneracy_2}
Consider the family \eqref{eq:polynomial_system_2}--\eqref{eq:linear_parameterization_2} and fix a point $x_*\in\R^n$. If there exists a parameter $\xi^0\in\Xi$ such that
\[
    Q_{\xi^0}(x_*)\ne0,
\]
then the set of parameters
\[
    \Sigma_{x_*}=\{\xi\in\Xi:Q_\xi(x_*)=0\}
\]
has zero Lebesgue measure in $\R^m$.
%is contained in a proper algebraic set and therefore has zero Lebesgue measure. Consequently, for almost all $\xi\in\Xi$ we have $Q_\xi(x_*)\ne0$. In particular, if for such $\xi$ we additionally have $x_*\in K_\xi$, then the nondegeneracy condition of Assumption \ref{ass:main_2} is satisfied.
\end{proposition}

It is proposed to give answers further in measure‑theoretic terms: “for almost all”, “up to a set of measure zero”, and the like. The reason is that for any somewhat complicated nonlinear system one can usually choose parameters in a certain degenerate way so that some conclusion fails. Declaring the parameters unknown, one cannot completely exclude degenerate cases (without adding a large number of complicated, hard‑to‑verify specific conditions); the best one can claim is that the chance of these “bad” cases is zero.

Thus, we formulate the main question of the paper as follows: suppose Assumption \ref{ass:main_2} hold for system \ref{eq:polynomial_system_2}; is it true that for $\mu$-almost every initial condition $x\in K$ the trajectory $t\mapsto\Phi^t(x)$ is affinely integrally exciting?



\section{Main Result}
\label{sec:main_result_2}


 The answer to the question of the previous section turns out to be positive.


For an affine hyperplane $H\subset\R^n$ set
\[
    U_\eps(H)=\{y\in K:\dist(y,H)<\eps\}.
\]


\begin{theorem}[Absence of sticking to hyperplanes and affine IPE]
\label{thm:main_nonsticking_2}
Let system \ref{eq:polynomial_system_2} satisfy Assumption \ref{ass:main_2}. Then there exists an invariant set $X_0\subset K$ of full measure, $\mu(X_0)=1$, such that for every $x\in X_0$ and every affine hyperplane $H\subset\R^n$ the fraction of time spent in the $\epsilon$-neighborhood of the hyperplane $U_\eps(H)$ tends to zero:
\begin{equation}\label{eq:main_nonsticking_2}
    \lim_{\eps\to0}\limsup_{T\to\infty}
    \frac1T\int_0^T\mathbf 1_{U_\eps(H)}(\Phi^t x)\dd t=0.
\end{equation}

\end{theorem}

\begin{proof}
    The proof of the theorem is given in the Appendix. For simplicity of verification, a reference to a Lean formalization of this proof is also provided there.
\end{proof}

\begin{corollary}
   Consequently, for every $x\in X_0$ the trajectory $t\mapsto\Phi^t(x)$ is affinely integrally exciting:
\begin{equation}\label{eq:main_ipe_2}
    \liminf_{T\to\infty}
    \frac1T\int_0^T \phi_z(t)\phi_z(t)^\top\dd t
    >0.
\end{equation}
\end{corollary}


\end{document}
